\documentclass[11pt]{article}

\usepackage[margin=1in]{geometry}
\usepackage{amsmath,amssymb}
\usepackage{graphicx}
\usepackage{booktabs}
\usepackage{hyperref}
\usepackage[caption=false]{subfig}

\title{Simulation-Based Inference under Model Misspecification\\with TabPFN}
\author{}
\date{}

\begin{document}
\maketitle

%% ============================================================
\section{Introduction}
%% ============================================================

Simulation-based inference (SBI) methods estimate posterior distributions
$p(\theta \mid x_{\mathrm{obs}})$ when the likelihood is intractable but a
simulator $p(x \mid \theta)$ is available.
A fundamental challenge arises when the simulator is \emph{misspecified}:
the true data-generating process differs from the assumed model, so the
simulated data distribution does not match the real observations.
Standard neural posterior estimation (NPE) methods can produce
overconfident or biased posteriors under such misspecification.

In this report we evaluate five approaches for posterior inference under
controlled misspecification across four benchmark tasks.
Our main contribution is a \emph{mixed-context} strategy
(NPE-PFN Mixed) that combines a small calibration dataset of
$(\theta, x)$ pairs from the true process with simulated data from the
misspecified model, leveraging TabPFN's in-context learning ability.

%% ============================================================
\section{Methods}
%% ============================================================

We compare five methods.  All NPE-PFN variants use TabPFN as the
amortised density estimator.

\begin{description}
  \item[NPE-PFN Misspec.]
    TabPFN conditioned only on simulated data from the misspecified
    model. No calibration data is used. This isolates the effect of
    misspecification.

  \item[NPE-PFN Calib.]
    TabPFN conditioned only on the calibration dataset
    $\{(\theta_i, x_i)\}_{i=1}^{n_{\mathrm{calib}}}$ drawn from the
    true process. No simulation data is used.

  \item[NPE-PFN Mixed (ours)]
    TabPFN with a mixed context: simulated rows are formatted as
    $(x_{\mathrm{sim}},\, \texttt{NaN})$ and calibration rows as
    $(x_{\mathrm{calib}},\, y_{\mathrm{calib}})$.
    At inference the query is $(\texttt{NaN},\, y_{\mathrm{obs}})$.
    This lets the model leverage both data sources simultaneously.

  \item[NPE (sbi)]
    A standard NPE baseline using a neural spline flow, trained on the
    calibration data via the \texttt{sbi} library.

  \item[NPE-PFN + FMPE]
    A two-stage method: (1)~a TabPFN $y$-corrector predicts $y$ from
    $(\theta, x)$, then (2)~FMPE is trained on synthetic
    $(\theta, \tilde y)$ pairs.  Skipped for SIR due to the
    computational cost of generating simulations through the ODE solver.
\end{description}

\subsection{Metrics}

We evaluate posterior quality with two complementary metrics.

\paragraph{C2ST (Classifier Two-Sample Test).}
A multi-layer perceptron is trained to distinguish approximate posterior
samples from reference samples using 5-fold cross-validation.
The reported accuracy lies in $[0.5, 1.0]$; a value of $0.5$
indicates identical distributions (chance level).
Lower is better.

\paragraph{MMD (Maximum Mean Discrepancy).}
The unbiased estimator of $\mathrm{MMD}^2$ with a Gaussian RBF kernel,
whose bandwidth is set via the median heuristic over pooled samples.
A value of $0$ indicates identical distributions.
Lower is better.

\subsection{Experimental setup}

For each task we draw $n_{\mathrm{sim}} = 5{,}000$ simulations from the
misspecified model ($10{,}000$ for SIR) and sweep the calibration set
size $n_{\mathrm{calib}} \in \{10, 50, 200, 1000\}$.
All results are averaged over 5 random seeds and 3 test observations
from the \texttt{sbibm} reference set.
Each method produces $n_{\mathrm{posterior}} = 1{,}000$ posterior
samples per observation.

%% ============================================================
\section{Benchmark Tasks}
%% ============================================================

% --------------------------------------------------
\subsection{Gaussian Mixture}
% --------------------------------------------------

The Gaussian Mixture task defines a two-component mixture model.
Parameters control the location factors and a mixture index.
The \textbf{one\_gaussian} misspecification replaces the mixture with a
single Gaussian using only the first component's location factor,
discarding the second component entirely:
\[
  x \sim \mathcal{N}\!\bigl(\mathrm{loc\_factor}_1 \cdot \theta,\;
  \sigma^2 I\bigr).
\]

\begin{figure}[htbp]
  \centering
  \includegraphics[width=\textwidth]{../results/gaussian_mixture_one_gaussian/sweep_avg.pdf}
  \caption{Gaussian Mixture (one\_gaussian misspec): C2ST and MMD vs.\
    calibration set size, averaged over seeds and observations.}
  \label{fig:gm-sweep}
\end{figure}

\begin{figure}[htbp]
  \centering
  \subfloat[Posterior pairplot]{%
    \includegraphics[width=0.48\textwidth]{../results/gaussian_mixture_one_gaussian/pairplots/ncalib50_seed42_obs1.pdf}%
    \label{fig:gm-pairplot}%
  }\hfill
  \subfloat[$y$-scatter diagnostic]{%
    \includegraphics[width=0.48\textwidth]{../results/gaussian_mixture_one_gaussian/y_scatter/ncalib50_seed42.pdf}%
    \label{fig:gm-yscatter}%
  }
  \caption{Gaussian Mixture: (a) posterior pairplot and (b) $y$-scatter
    diagnostic at $n_{\mathrm{calib}}=50$, seed 42, observation 1.}
\end{figure}

% --------------------------------------------------
\subsection{Two Moons}
% --------------------------------------------------

The Two Moons task generates data from a crescent-shaped distribution.
The \textbf{heavy\_tail\_radius} misspecification replaces the Gaussian
radial noise $r \sim \mathcal{N}(0.1, 0.01)$ with a Student-$t$
distribution $r \sim t_\nu(0.1, 0.01)$ with $\nu = 2$ degrees of
freedom, introducing heavier tails.

\begin{figure}[htbp]
  \centering
  \includegraphics[width=\textwidth]{../results/two_moons_heavy_tail_radius/sweep_avg.pdf}
  \caption{Two Moons (heavy\_tail\_radius misspec): C2ST and MMD vs.\
    calibration set size.}
  \label{fig:tm-sweep}
\end{figure}

\begin{figure}[htbp]
  \centering
  \subfloat[Posterior pairplot]{%
    \includegraphics[width=0.48\textwidth]{../results/two_moons_heavy_tail_radius/pairplots/ncalib50_seed42_obs1.pdf}%
    \label{fig:tm-pairplot}%
  }\hfill
  \subfloat[$y$-scatter diagnostic]{%
    \includegraphics[width=0.48\textwidth]{../results/two_moons_heavy_tail_radius/y_scatter/ncalib50_seed42.pdf}%
    \label{fig:tm-yscatter}%
  }
  \caption{Two Moons: (a) posterior pairplot and (b) $y$-scatter
    diagnostic at $n_{\mathrm{calib}}=50$, seed 42, observation 1.}
\end{figure}

% --------------------------------------------------
\subsection{SLCP}
% --------------------------------------------------

The SLCP (Simple Likelihood Complex Posterior) task produces 8-dimensional
observations from a 5-dimensional parameter space.
The \textbf{additive\_noise} misspecification adds i.i.d.\ Gaussian noise
to each output dimension:
\[
  x_{\mathrm{misspec}} = x_{\mathrm{true}} +
  \mathcal{N}(0,\, \sigma_{\mathrm{noise}}^2 I),
  \quad \sigma_{\mathrm{noise}} = 0.5.
\]

\begin{figure}[htbp]
  \centering
  \includegraphics[width=\textwidth]{../results/slcp_additive_noise/sweep_avg.pdf}
  \caption{SLCP (additive\_noise misspec): C2ST and MMD vs.\
    calibration set size.}
  \label{fig:slcp-sweep}
\end{figure}

\begin{figure}[htbp]
  \centering
  \includegraphics[width=0.55\textwidth]{../results/slcp_additive_noise/pairplots/ncalib50_seed42_obs1.pdf}
  \caption{SLCP: posterior pairplot at $n_{\mathrm{calib}}=50$, seed 42,
    observation 1.}
  \label{fig:slcp-pairplot}
\end{figure}

% --------------------------------------------------
\subsection{SIR}
% --------------------------------------------------

The SIR task models an epidemic via a susceptible--infected--recovered
ODE with parameters for the infection rate $\beta$, recovery rate
$\gamma$, and initial infected fraction.
Observations are daily infection counts with Binomial observation noise.

The \textbf{weekend\_delay} misspecification (Ward et al., 2022)
simulates reporting delays: a fraction of Saturday and Sunday infection
counts is shifted to the following Monday.
The delay fraction is set to $0.05$.
This produces characteristic underreporting on weekends with a Monday
catch-up spike.

The NPE-PFN+FMPE method is skipped for this task because generating
$10{,}000$ synthetic $y$-values through the ODE solver is
prohibitively expensive.

\begin{figure}[htbp]
  \centering
  \includegraphics[width=\textwidth]{../results/sir_weekend_delay/sweep_avg.pdf}
  \caption{SIR (weekend\_delay misspec): C2ST and MMD vs.\
    calibration set size.}
  \label{fig:sir-sweep}
\end{figure}

\begin{figure}[htbp]
  \centering
  \includegraphics[width=0.55\textwidth]{../results/sir_weekend_delay/pairplots/ncalib50_seed42_obs1.pdf}
  \caption{SIR: posterior pairplot at $n_{\mathrm{calib}}=50$, seed 42,
    observation 1.}
  \label{fig:sir-pairplot}
\end{figure}

%% ============================================================
\section{Discussion}
%% ============================================================

Across all tasks, the NPE-PFN Misspec.\ baseline degrades substantially
when the simulator departs from the true data-generating process,
confirming that standard amortised inference is sensitive to
misspecification.
The calibration-only method (NPE-PFN Calib.) performs well when
sufficient calibration data is available but is data-hungry.

The mixed-context approach (NPE-PFN Mixed) consistently achieves
competitive or superior performance, particularly in the low-data regime
($n_{\mathrm{calib}} \leq 50$), by combining the inductive bias from
simulations with the distributional correction from calibration data.
This confirms the value of leveraging TabPFN's in-context learning to
fuse heterogeneous data sources.

The NPE (sbi) baseline, trained only on calibration data with a neural
spline flow, generally requires more data to match the TabPFN-based
methods.  The two-stage NPE-PFN+FMPE approach offers an alternative
correction strategy but adds computational overhead and complexity.

The SIR task introduces a more realistic misspecification
(weekend reporting delays) with higher-dimensional observations.
All methods struggle on this task, with C2ST values near $1.0$,
indicating that the high-dimensional time-series observation space
poses a significant challenge for posterior estimation under
misspecification.

\end{document}
